\chapter*{AI Usage}

I would like to transparently clarify how I used AI in this thesis project.

\paragraph{Research} At the beginning of the project, while exploring potential design ideas for Boogie map semantics
(see \Cref{sec:support-boogie-maps})
I used Gemini 3 and ChatGPT based on GPT-5.3 to generate possible approaches.
These attempts were less successful, as they mostly suggested common
association-list-based approaches (see \Cref{sec:assoc-list}) or context/heap-model ideas (see \Cref{sec:contex-model}).
Such approaches are useful for modeling classical dictionaries in mainstream languages,
but not for capturing the functional nature of Boogie maps.

\paragraph{Proof Engineering} Most of the proofs written in Isabelle/HOL were written by hand, with little or no AI assistance.
One extensionality proof was guided by Gemini 3. It produced an incorrect proof, which I then corrected.
Another area in which AI was helpful was proving lifted versions of definitions and lemmas, as I initially lacked expertise with the lifting package.
Here, Gemini 3 was helpful as well.
Otherwise, Gemini and ChatGPT were used mainly to explain errors or suggest references for alternative proof strategies.

\paragraph{Programming} For programming and adjusting the proof generation module,
AI was used to debug a few difficult cases, explain existing code, and suggest small code snippets.
For this, I used Junie by JetBrains and Claude Agent by Anthropic, both within JetBrains Rider.

\paragraph{\LaTeX} I used OpenAI Prism as my main editor to create this thesis in \LaTeX.
I used it to convert Markdown to \LaTeX, create TikZ figures from visuals I had created in a slide deck,
generate tables, and fix \LaTeX{} errors.

\paragraph{Writing} I used AI during the writing of this thesis.
My workflow was to write a draft by hand and then have it proofread and reformulated by an LLM.
I primarily used Gemini 3 and OpenAI Prism.
I used them to correct typographical and grammatical errors and to improve sentence structure
and academic tone.
